\documentclass[11pt]{article}
\usepackage[a4paper,margin=2.6cm]{geometry}
\usepackage{amsmath,amssymb,amsthm,mathtools}
\usepackage{booktabs}
\usepackage[colorlinks=true,linkcolor=blue,citecolor=blue,urlcolor=blue]{hyperref}
\usepackage{microtype}

\newtheorem{theorem}{Theorem}
\newtheorem{proposition}[theorem]{Proposition}
\newtheorem{lemma}[theorem]{Lemma}
\newtheorem{corollary}[theorem]{Corollary}
\newtheorem{fact}[theorem]{Fact}
\theoremstyle{definition}
\newtheorem{definition}[theorem]{Definition}
\newtheorem{remark}[theorem]{Remark}

\title{\texorpdfstring{The Alternating Series $\sum_{n\ge1}(-1)^n n/p_n$: Exact Computations, Reductions, and the Exact Obstruction (Erd\H{o}s Problem EP-15)}{The Alternating Series of Erdos Problem EP-15: Exact Computations, Reductions, and the Exact Obstruction}}
\author{ox-alpha (orchestrated multi-agent analysis)}
\date{August 26, 2026}

\begin{document}
\maketitle

\begin{abstract}
\noindent
We consider the alternating series $\sum_{n\ge1}(-1)^n n/p_n$, where $p_n$ denotes the $n$-th prime (Erd\H{o}s Problem EP-15).
\emph{(i)}~The problem is open unconditionally; no proof or disproof of convergence is known, and we do not claim one.
\emph{(ii)}~Convergence holds under the Quantitative Hardy--Littlewood prime tuples conjecture (Conjecture~1.3 and Theorem~1.4 of Tao \cite{Ta23}; published in \cite{Ta24}).
\emph{(iii)}~Exact fixed-point computations (scale $10^{60}$, rigorous rounding control) verify convergence-consistency through $N=10^7$.
\emph{(iv)}~The classical alternating-series machinery is provably inapplicable: bounded gaps kill monotonicity, GPY positive-density small gaps force total variation $\gg\ln\ln N$, killing the Dirichlet criterion.
\emph{(v)}~A counterexample ``world'' $W(\theta,\rho)$ reproducing every published unconditional prime statistic (now including the full second-order Cipolla term) while its paired series diverges shows that current unconditional technology cannot decide the problem.
\emph{(vi)}~The exact missing implication --- uniform parity-resolved gap statistics equivalent to parity equidistribution of $\pi(m)$ in log-length intervals --- is isolated; it is supplied today only by Hardy--Littlewood strength input.
\end{abstract}

\section{Problem and verdict}\label{sec:problem}

\begin{quote}
\textbf{Question (EP-15, verbatim).} Is it true that
\[
S \;=\; \sum_{n=1}^{\infty} (-1)^n \frac{n}{p_n}
\]
converges, where $p_n$ denotes the $n$-th prime?
\end{quote}

\begin{center}
\fbox{\parbox{0.94\linewidth}{
\textbf{VERDICT.}
\begin{itemize}\itemsep1pt
\item \textbf{UNCONDITIONALLY OPEN.} No unconditional proof or disproof of convergence is known. This investigation does \emph{not} solve the problem and does not claim to.
\item \textbf{CONDITIONAL [PROVEN under HL-strong]:} convergence holds assuming the Quantitative Hardy--Littlewood prime tuples conjecture (Conjecture~1.3 of Tao \cite{Ta23}: tuples of size $k\le(\log\log x)^5$ in $[0,\log^2 x]$ with power-savings error $Cx^{1-\varepsilon}$), by Theorem~1.4 of \cite{Ta23}.
\item \textbf{NUMERICAL: CONVERGENCE-CONSISTENT through $N=10^7$} at exact fixed-point precision (Section~\ref{sec:num}): steadily shrinking half-decade increments, no drift signal, paired partial sums bounded after the transient first pair.
\item The problem is \textbf{not solved and not disproved} here. Delivered instead: proofs that the classical criteria are inapplicable (Sections~\ref{sec:r1}, \ref{sec:r2}), an exact reduction to an isolated parity-resolved open subproblem (Sections~\ref{sec:r3}, \ref{sec:r4}), a counterexample world showing no current unconditional theorem can decide the question (Section~\ref{sec:r5}), and integration of the numerical evidence (Section~\ref{sec:num}).
\end{itemize}}}
\end{center}

\section{Notation and the three exact identities}\label{sec:not}

Throughout,
\[
a_n:=\frac{n}{p_n},\qquad g_n:=p_{n+1}-p_n,\qquad S_N:=\sum_{n\le N}(-1)^n a_n,\qquad L_k:=\ln(2k),\qquad \ell_k:=\ln\ln(2k),\qquad \Lambda_k:=L_k+\ell_k .
\]
The series starts with $(-1)^1a_1=-1/2$, so pairing consecutive terms $(1,2),(3,4),\dots$ gives $S_{2K}=\sum_{k\le K}(a_{2k}-a_{2k-1})$ exactly.

The three exact identities on which everything rests are:

\begin{equation}\label{eq:mono}
a_{n+1}<a_n \;\Longleftrightarrow\; \frac{n+1}{p_{n+1}}<\frac{n}{p_n}
\;\Longleftrightarrow\; (n{+}1)p_n<np_{n+1}
\;\Longleftrightarrow\; n\,g_n>p_n ,
\end{equation}
and symmetrically $a_{n+1}>a_n\iff ng_n<p_n$; upward steps of $(a_n)$ occur exactly where the gap is below the local mean-gap scale $p_n/n$.

\begin{equation}\label{eq:step}
a_{n+1}-a_n=\frac{(n{+}1)p_n-np_{n+1}}{p_np_{n+1}}
=\frac{np_n+p_n-n(p_n+g_n)}{p_np_{n+1}}
=\frac{p_n-n\,g_n}{p_n\,p_{n+1}},
\end{equation}
with sign convention positive step $\iff p_n>ng_n$, consistent with \eqref{eq:mono}.

\begin{equation}\label{eq:pair}
b_k:=a_{2k}-a_{2k-1}
=\frac{2k\,p_{2k-1}-(2k-1)p_{2k}}{p_{2k-1}p_{2k}}
=\frac{p_{2k-1}-(2k-1)g_{2k-1}}{p_{2k-1}p_{2k}},
\qquad
b_k>0\iff g_{2k-1}<\frac{p_{2k-1}}{2k-1},
\end{equation}
where the numerator rewrite uses $p_{2k}=p_{2k-1}+g_{2k-1}$ and $2kp_{2k-1}-(2k-1)(p_{2k-1}+g_{2k-1})=p_{2k-1}-(2k-1)g_{2k-1}$.
Sanity checks (exact rational arithmetic against an independent sieve, all $k\le 50{,}000$): both closed forms in \eqref{eq:pair} equal $a_{2k}-a_{2k-1}$ everywhere tested; $b_1=a_2-a_1=\tfrac23-\tfrac12=\tfrac16=(2-1)/6$, matching $\max_K|B_K|=1/6$ attained at $K=1$ in the data of Section~\ref{sec:num}.

\section{R1: the Leibniz test is provably inapplicable}\label{sec:r1}

Leibniz requires $(a_n)$ eventually non-increasing. It fails in both directions.

\begin{fact}[upward steps infinitely often]\label{fac:up}
By Maynard's strengthening of Zhang \cite{Zh14,Ma15,Po14}, $\liminf_n g_n\le 246$, so $g_n\le 246$ infinitely often. By the prime number theorem $p_n/n\sim\ln n\to\infty$; choose $n_0$ with $p_n/n>246$ for $n\ge n_0$. Infinitely many such indices lie beyond $n_0$, and there $ng_n\le246n<p_n$, i.e.\ $a_{n+1}>a_n$ by \eqref{eq:mono}.
\end{fact}

\begin{fact}[downward steps infinitely often]\label{fac:down}
Westzynthius \cite{We31} proved $\limsup g_n/\ln p_n=+\infty$; hence $g_n>\tfrac32\ln p_n$ infinitely often. Since $p_n\ge n$, $\ln p_n\ge\ln n$, so $g_n>\tfrac32\ln n$ infinitely often; while PNT gives $p_n/n=(1+o(1))\ln n\le\tfrac54\ln n$ for $n\ge n_1$. Along infinitely many indices both hold: $ng_n>p_n$, i.e.\ $a_{n+1}<a_n$.
\end{fact}

\begin{corollary}
$(a_n)$ is eventually monotone in neither direction; the Leibniz test cannot be formulated, unconditionally. \textnormal{[NUMERICAL cross-check: $42.13\%$ of steps $n\le10^7$ are downward.]}
\end{corollary}

\section{R2: divergence of total variation and failure of Dirichlet--BV}\label{sec:r2}

Dirichlet's test in BV form: if $a_n\to0$ and $(a_n)$ has bounded total variation then $\sum(-1)^na_n$ converges (Abel summation). We show total variation diverges unconditionally.

Input (unconditional, cited): \emph{GPY positive-proportion small gaps} \cite[Theorem~1]{GPY11}: for every fixed $\eta>0$,
\[
P(x,\eta):=\frac{1}{\pi(x)}\#\{p_n\le x:\ g_n\le\eta\ln p_n\}\gg_\eta 1 .
\]

\begin{lemma}[Window Lemma]\label{lem:window}
Let $A\subseteq\mathbb N$ have lower density $\delta>0$, i.e.\ $\#(A\cap[1,N])\ge\delta N$ for all $N\ge N_0$. Then
$\sum_{n\in A,\,n\le N}\frac{1}{n\ln n}\gg_\delta \ln\ln N$.
\end{lemma}
\begin{proof}
Let $r:=4/\delta\,(>2)$ and $x_j:=x_0r^j$ with $x_0:=\max(2N_0,r)$. Each window $W_j:=(x_j/r,\,x_j]$ satisfies
$\#\,(A\cap W_j)=A(x_j)-A(\lfloor x_j/r\rfloor)\ \ge\ \delta x_j-x_j/r\ \ge\ (\delta/2)x_j$
(using only the total lower density, no blockwise hypothesis). The windows are pairwise disjoint, lie in $[1,N]$ whenever $x_j\le N$, and
\[
\sum_{n\in A\cap W_j}\frac{1}{n\ln n}\ \ge\ \#\,(A\cap W_j)\cdot\frac{1}{x_j\ln x_j}\ \ge\ \frac{\delta}{2\ln x_j}
=\frac{\delta}{2(j\ln r+\ln x_0)},
\]
and summing over $j\le\log_r(N/x_0)$ gives $\gg_\delta\ln\ln N$.
\end{proof}

\begin{proposition}\label{prop:tv}
Fix $\eta=\tfrac14$ and $A=\{n:\ g_n\le\eta\ln p_n\}$, of lower density $\delta=\delta(\eta)>0$. Then for some absolute $c>0$ and all large $N$,
\[
TV(N):=\sum_{n<N}|a_{n+1}-a_n|\ \ge\ \sum_{\substack{n\le N\\ n\in A}}\frac{c}{n\ln n}\ \gg\ \ln\ln N .
\]
\end{proposition}
\begin{proof}
Step 1 (pointwise). For $n\in A$, by \eqref{eq:step},
\[
|a_{n+1}-a_n|=\frac{p_n-ng_n}{p_np_{n+1}}\ \ge\ \frac{p_n-\eta\,n\ln p_n}{p_np_{n+1}}
=\frac{1}{p_{n+1}}\Bigl(1-\eta\,\frac{n\ln p_n}{p_n}\Bigr).
\]
Since $p_n/n\sim\ln n$, we have $n\ln p_n/p_n\to1$; with $\eta=1/4$ the bracket is $\ge1/2$ for $n\ge n_2$, and $p_{n+1}\sim n\ln n$, so $|a_{n+1}-a_n|\ge c/(n\ln n)$ on $A$.
Step 2. Combine with Lemma~\ref{lem:window}.
\end{proof}

\begin{corollary}
$TV(N)\to\infty$ unconditionally (empirically like $0.55\ln\ln N$, Section~\ref{sec:num}); the BV hypothesis fails, so Dirichlet--BV provably cannot decide EP-15. Of course $TV\to\infty$ does not imply divergence of the series (cf.\ $\sum(-1)^n/\sqrt{n}$); it only kills this route.
\end{corollary}

\section{R3: smooth/fluctuation decomposition and the parity-blind residual}\label{sec:r3}

\textbf{Cipolla expansion} \cite{Ci02}. With $L=\ln n$, $\ell=\ln\ln n$,
\[
p_n=n\Bigl(L+\ell-1+\frac{\ell-2}{L}-\frac{\ell^2-6\ell+11}{2L^2}+O\Bigl(\frac{\ell^3}{L^3}\Bigr)\Bigr),
\]
and we fix the truncation
\[
E(n):=n\Bigl(L+\ell-1+\frac{\ell-2}{L}-\frac{\ell^2-6\ell+11}{2L^2}\Bigr),\qquad f(n):=\frac{n}{E(n)},
\]
so that $|p_n-E(n)|\ll n\,\ell^3/L^3$ (omitted-term scale) and moreover $|p_n-E(n)|\ll ne^{-c\sqrt{\ln n}}$ by inversion of the de la Vall\'ee Poussin error $\pi(x)=\operatorname{li}(x)+O(xe^{-c\sqrt{\ln x}})$ \cite{vK01,DVLP}.

\begin{lemma}[$f$ is eventually strictly decreasing]\label{lem:f}
Writing $u=\ln n$ and $G(u)=L+\ell-1+\frac{\ell-2}{L}-\frac{\ell^2-6\ell+11}{2L^2}$ (so $E(n)=nG(u)$, $f(n)=1/G(u)$),
\[
f'(n)=-\frac{G'(u)/n}{G(u)^2},\qquad
G'(u)=1+\frac1u+\frac{3-\ln u}{u^2}+\frac{(\ln u)^2-7\ln u+14}{u^3},
\]
an identity verified symbolically by CAS. Since $(\ln u)^2-7\ln u+14$ has negative discriminant ($49-56<0$) it is positive for all $u$, so $G'(u)\ge\tfrac12$ for $u\ge e^3$, giving $-\frac{2}{nG(u)^2}\le f'(n)\le-\frac{1}{2nG(u)^2}$ there. Hence $f'<0$ eventually and $f(n)\asymp1/\ln n\to0$; by Leibniz, $\sum(-1)^nf(n)$ converges with remainder $\le f(N{+}1)\ll1/\ln N$.
\end{lemma}

\begin{lemma}[exact residual identity]\label{lem:delta}
With
\[
\delta_n:=a_n-f(n)=\frac{n}{p_n}-\frac{n}{E(n)}
=-\,(p_n-E(n))\cdot\frac{n}{p_nE(n)},
\]
and $p_n,E(n)\asymp n\ln n$ (PNT),
\begin{equation}\label{eq:env}
|\delta_n|\ \lesssim\ \frac{|p_n-E(n)|}{n\ln^2 n}\ \ll\ \min\Bigl(\frac{\ell^3}{L^3},\ e^{-c\sqrt{\ln n}}\Bigr).
\end{equation}
Consequently the original series converges $\iff \sum(-1)^n\delta_n$ converges.
\end{lemma}

\paragraph{Why the residual is stuck unconditionally.}
The envelope \eqref{eq:env} decays slower than any power and is non-summable; worse, $\delta_n$ inherits the erratic sign of $p_n-E(n)$ and no monotonicity or parity structure is known for it. All zero-free-region technology bounds the envelope $|p_n-E(n)|$ and is \emph{parity-blind}: it provides no information on $(-1)^n$-correlations of the fluctuations.

\paragraph{What RH-grade input would (and would not) give [CONDITIONAL: RH].}
Under RH, von Koch \cite{vK01}: $\pi(x)=\operatorname{li}(x)+O(\sqrt{x}\ln x)$; inversion yields $p_n=\operatorname{li}^{-1}(n)+O(\sqrt n\ln^2n)$, whence $|\delta_n|=O(n^{-1/2+o(1)})$. But envelope decay alone does not suffice: the comparison sequence $\delta'_n=(-1)^n/(\sqrt n\ln n)$ meets the same envelope while $\sum(-1)^n\delta'_n=\sum1/(\sqrt n\ln n)$ diverges. The Dirichlet-eta phenomenon ($\sum(-1)^nn^{-s+i\tau}$ convergent for $\Re s>0$) relies on rigid monomial structure; interchanging limits over zeros (explicit-formula expansions followed by alternating summation) is not available in any published form. \emph{No published unconditional or RH-only proof exists; the published conditional route is HL-strength (Tao).} Marked honestly: Lemmas~\ref{lem:f}--\ref{lem:delta} and \eqref{eq:env} are theorems; ``RH reopens an eta-type argument'' is plausible-but-unpublished.

\paragraph{Cross-reference to the published route.}
Tao \cite{Ta23}, recording an unpublished observation of Said \cite{MO}, proves the asymptotic relation
\begin{equation}\label{eq:said}
\sum_{n\le x}\frac{(-1)^nn}{p_n}=\frac12\sum_{2\le m\le x/\log x}\frac{(-1)^{\pi(m)}}{m\log m}+C+o(1),
\end{equation}
reducing EP-15 to cancellation in $(-1)^{\pi(m)}$; under Conjecture~1.3 of \cite{Ta23} he obtains $\bigl|\sum_{n\le x}(-1)^{\pi(n)}\bigr|\ll x/(\log\log x)^{1.1}$ [CONDITIONAL], via Gallagher-type Poisson statistics \cite{Ga76} quantified by Kuperberg \cite{Ku} and the Banks--Ford--Tao random sifted model \cite{BFT}.

\section{R4: exact pair reduction}\label{sec:r4}

Define the pair terms $b_k=a_{2k}-a_{2k-1}$. Then $S_{2K}=\sum_{k\le K}b_k$ exactly, and $S_{2K+1}=S_{2K}-a_{2K+1}$ with $a_{2K+1}\to0$; therefore
\begin{equation}\label{eq:equiv}
\boxed{\ \sum_{n\ge1}(-1)^n\frac np_n\ \text{converges}\ \iff\ \sum_{k}b_k\ \text{converges}.\ }
\end{equation}

\begin{equation}\label{eq:pairbox}
\boxed{\ b_k=\frac{p_{2k-1}-(2k-1)g_{2k-1}}{p_{2k-1}p_{2k}},\qquad
b_k>0\iff g_{2k-1}<\frac{p_{2k-1}}{2k-1}\ }
\end{equation}
with the sign convention of Section~\ref{sec:not}: positivity means the even-indexed term is larger, i.e.\ the odd-indexed gap is \emph{below} the threshold $p_{2k-1}/(2k-1)\approx\ln(2k)$.

\begin{remark}[sign erratum vs.\ the original tasking sheet]
The sheet asserted ``$b_k>0\iff g_{2k-1}>p_{2k-1}/(2k-1)$'' and offered the numerator $(2k-1)g-p$; both are sign-inverted. The correct chain, verified above, is $2kp_{2k-1}-(2k-1)p_{2k}=p_{2k-1}-(2k-1)g_{2k-1}$. Intuition confirms: a smaller-than-average gap before $p_{2k}$ makes $n/p_n$ jump up across the pair. Qualitative conclusions unaffected.
\end{remark}

At leading order, with $p_{2k-1},p_{2k}\asymp2kL_k$,
\begin{equation}\label{eq:lead}
b_k\asymp\frac{L_k-g_{2k-1}}{2kL_k^2}=\frac{1-g_{2k-1}/L_k}{2kL_k},
\end{equation}
so $|b_k|$ has typical size $\asymp1/(k\ln k)$ --- matching the observed $TV(N)\sim0.55\ln\ln N$ --- while its \emph{sign} is governed by whether the odd-indexed normalized gap falls below $1$. Convergence of $\sum_kb_k$ is thus a statement about \emph{parity-resolved statistics of normalized gaps}: persistent odd/even bias creates a one-signed drift $\sim\frac{|1-\theta|}{2}\ln\ln K$. This is the isolated open subproblem, equivalent via \eqref{eq:said} to parity-equidistribution of $\pi(m)$ in log-length intervals.

\section{R5: counterexample world \texorpdfstring{$W(\theta,\rho)$}{W(theta,rho)}}\label{sec:r5}

We construct a hypothetical arithmetic world --- a strictly increasing integer sequence $(\tilde p_j)$ playing the role of primes --- exhibiting persistent odd/even gap bias yet satisfying every published unconditional property of the primes at current strength. Items marked (P) are proved here; compatibility items marked [HEURISTIC] are enforced by construction up to standard balancing arguments.

\begin{definition}[$W(\theta,\rho)$]\label{def:world}
Fix $\theta\in(0,2)$, $\theta\ne1$, $\rho\in[0,1)$; recall $L_k=\ln(2k)$, $\ell_k=\ln\ln(2k)$, $\Lambda_k=L_k+\ell_k$.
\begin{enumerate}\itemsep2pt
\item \emph{(Bulk, P)} For $k\notin\mathcal A$ (below), set $\tilde g_{2k-1}=\theta\Lambda_k(1+\xi_k)$, $\tilde g_{2k}=(2-\theta)\Lambda_k(1+\xi'_k)$ with balanced deterministic signs $\xi,\xi'=\pm o(1)$; pair totals are exactly preserved: $\tilde g_{2k-1}+\tilde g_{2k}=2\Lambda_k=2(\ln(2k)+\ln\ln(2k))$ when $\xi=\xi'=0$. (The $\ell_k$ layer is what makes the world reproduce the full Cipolla main terms; see Proposition~\ref{prop:tel}.)
\item \emph{(Sprinkle, P)} On a set $\mathcal A$ of lower density $\rho$ among odd positions, shrink the odd gap into $[2,246]$, and redistribute the borrowed mass $\Delta_k=\theta\Lambda_k-O(1)$ onto odd compensators at positions $k'$ with $k<k'\le k+k^{1/2}$ (same parity!), so each event contributes to $\tilde S_{2K}$ only its mismatch: $\bigl|\frac{1}{2k\Lambda_k}-\frac{1}{2k'\Lambda_{k'}}\bigr|\lesssim\frac{k'-k}{k^2\Lambda_k}\lesssim\frac{1}{k^{3/2}\Lambda_k}$.
\item \emph{(Pulses, P)} Enlarging $\tilde g_{2k}$ by $\Delta$ and shrinking $\tilde g_{2k+2}$ by $\Delta$ keeps all later cumulatives unchanged except $\tilde p_{2k+1}$, lifted by $\le O(\Lambda_k)$.
\end{enumerate}
\end{definition}

\subsection{Pair-drift divergence; convergent sprinkle cost}

\begin{proposition}[drift law unchanged under the $\Lambda_k$ pattern]\label{prop:drift}
$\displaystyle \sum_{k\le K}\tilde b_k=\frac{1-\theta}{2}\ln\ln K+O(1)$, divergent for every $\theta\ne1$.
\end{proposition}
\begin{proof}
With $\tilde g_{2k-1}=\theta\Lambda_k$ and the uniform baseline of Proposition~\ref{prop:tel} below, $\tilde p_m=m(\ln m+\ln\ln m-1)+O(m/\ln m)$, so $\tilde p_{2k-1}=(2k-1)(L_k+\ell_k-1)+O(k/\ln k)$ (off-by-one costs $O(1)$ in each factor) and $\tilde p_{2k}=2k\Lambda_k+O(k/\ln k)$. Substituting into the exact identity \eqref{eq:pair} applied in-world,
\[
\tilde b_k=\frac{\tilde p_{2k-1}-(2k-1)\theta\Lambda_k}{\tilde p_{2k-1}\tilde p_{2k}}
=\frac{(2k-1)\bigl[(1-\theta)\Lambda_k-1\bigr]+O(k/\ln k)}{(2k)^2\Lambda_k^2(1+o(1))},
\]
which decomposes into three series:
\[
\underbrace{\frac{(1-\theta)\Lambda_k}{2k\Lambda_k^2}\sim\frac{1-\theta}{2kL_k}}_{\text{parity drift}},\qquad
\underbrace{-\frac{1}{2k\Lambda_k^2}}_{\text{systematic }-1\text{ term}},\qquad
\underbrace{O\Bigl(\frac{1}{k\Lambda_k^2\ln k}\Bigr)}_{\text{baseline error}} .
\]
The first sums like $\frac{1-\theta}{2}\sum_{k\le K}\frac1{k\ln k}\sim\frac{1-\theta}{2}\ln\ln K$ --- divergent for every $\theta<1$ (opposite sign for $\theta>1$); the second and third are $O(\sum1/(k\ln^2k))=O(1)$: convergent.
\end{proof}

Sprinkled events: a direct bounded-gap replacement contributes $\asymp\frac{\Lambda_k}{2k\Lambda_k^2}=\frac{1}{2k\Lambda_k}$, its same-parity compensator $-\frac{1}{2k'\Lambda_{k'}}$; net per event $\lesssim\frac{1}{k^{3/2}\ln k}$, and even granting all $\rho K$ events the trivial bound, $\sum_{k\ge2}\frac{1}{k^{3/2}\ln k}<\infty$. Pulse effects act on baselines only, $\partial\tilde b/\partial p\asymp1/p^2$, contributing $\lesssim\sum_k\frac{\Lambda_k}{k^2\Lambda_k^2}=\sum\frac{1}{k^2\ln k}<\infty$. Total sprinkle+pulse contribution: an absolutely convergent series of order $\sum O(1/(k^{3/2}\ln k))$ (same convergent class as the announced $\sum O(1/(k\ln^2k))$). Hence $W(\theta,\rho)$ has divergent $\sum_k\tilde b_k$ for every $\theta\ne1$, independent of $\rho$.

\begin{remark}[erratum vs.\ the executed probe script]
The script \texttt{ep15\_partial\_sums.py} (adversarial block) accumulated $(\theta-1)\sum_{k\le K}\frac{1}{2kL_k^2}$ --- denominator $L_k^2$ without the numerator factor $L_k$ from \eqref{eq:lead} --- whose total \emph{converges}; its printed values ($-0.722145,-0.144429,-0.014443$ at $K=10^6$ for $\theta=0.5,0.9,0.99$) therefore do not themselves exhibit $\ln\ln$-divergence, and its stated per-pair form $\frac{g-L_k}{2kL_k^2}$ carries the opposite sign to \eqref{eq:lead}. Corrected model drift at $K=10^6$: $(1-\theta)\sum_{k\le10^6}\frac{1}{2k\ln(2k)}=(1-\theta)(1.997918\dots)$, i.e.\ $+0.998959,+0.199792,+0.019979$ (sub-average odd gaps push $S_{2K}$ \emph{up}). An inline erratum comment has been inserted in the script; the qualitative conclusion --- persistent parity bias forces divergence --- survives and is proved above.
\end{remark}

\subsection{Telescoping: the world reproduces the full Cipolla main terms}

\begin{proposition}\label{prop:tel}
In $W(\theta,\rho)$,
\[
\boxed{\ \tilde p_N=N\bigl(\ln N+\ln\ln N-1\bigr)+O\Bigl(\frac{N}{\ln N}\Bigr).\ }
\]
That is, the world reproduces the full Cipolla main terms through the constant $-1$, with honestly stated residual $O(N/\ln N)=o(N)$.
\end{proposition}
\begin{proof}
Bulk pair totals are $2\Lambda_k$; redistribution moves mass only within odd positions across $k^{1/2}$-windows, so odd-position cumulative mass is conserved up to window-boundary leakage $\ll\rho K^{1/2}\ln K$ at $2K$; pulses have zero net effect on even cumulatives. Summing over pairs, the $\theta$-dependence cancels exactly in even-indexed cumulatives ($\theta\Sigma_\Lambda+(2-\theta)\Sigma_\Lambda=2\Sigma_\Lambda$):
\[
\tilde p_{2K}=2+\sum_{k\le K}(\tilde g_{2k-1}+\tilde g_{2k})=2+2\sum_{k\le K}\Lambda_k+O\bigl(\rho K^{1/2}\ln K\bigr).
\]
Euler--Maclaurin bookkeeping: $\sum_{k\le K}\ln(2k)=K\ln(2K)-K+O(\ln K)$ (Stirling), and, using the antiderivative $\frac{\mathrm d}{\mathrm dx}\bigl[x\ln\ln x-\operatorname{li}(x)\bigr]=\ln\ln x$ together with $\operatorname{li}(x)=\frac{x}{\ln x}(1+O(\frac1{\ln x}))$,
\[
\sum_{k\le K}\ell_k=\sum_{k\le K}\ln\ln(2k)=K\ln\ln(2K)-\tfrac12\operatorname{li}(2K)+O(\ln\ln K)
=K\ln\ln(2K)-\frac{K}{\ln(2K)}+O\Bigl(\frac{K}{\ln^2K}\Bigr).
\]
Therefore
$\tilde p_{2K}=2K\bigl(\ln(2K)+\ln\ln(2K)-1\bigr)+O\bigl(K/\ln K\bigr)$;
since one further increment costs $O(\ln N)$ and pulse lifts are $O(\ln N)$, the claim holds uniformly in $N$.
\end{proof}

Consequently $\tilde\psi(x),\tilde\vartheta(x)$ match the true-prime PNT scale with identical leading and second-order structure up to $o(N)$, and $\tilde\pi(x)/\pi(x)\to1$ at PNT scale. [HEURISTIC: enforcing the full de la Vall\'ee Poussin \emph{error} inside $W$ requires a standard discrepancy-balancing choice of the $\xi$'s; the main-term statement just given is proved.]

\subsection{Compatibility with published unconditional theorems}

\begin{itemize}\itemsep2pt
\item \emph{Bounded gaps i.o.\ and at density:} built in (sprinkle, size $\le246$, density $\rho>0$) $\to$ Zhang/Maynard/Polymath/Pintz-type statements satisfiable. (P)
\item \emph{Upper bounds on gaps:} bulk gaps $\le2\Lambda_k(1+\xi_{\max})+O(1)=O(\ln N)$ (second-order Cipolla scale included), $\ll p_N^{0.525}$ $\to$ BHP exponent compatible; PNT matched through the $\ln\ln N$ term. (P)
\item \emph{Large-gap limsup lower bounds} (Westzynthius/Erd\H{o}s--Rankin/Ford--Green--Konyagin--Maynard--Tao): insert Rankin-scale gaps at super-exponentially sparse positions $n_j$ (both parities; even ones pulse-compensated). Cost to $\sum_k\tilde b_k$ per event $\lesssim\Delta_j/(n_j\ln n_j)$ (odd), resp.\ baseline class $\lesssim\Delta_j/(n_j^2\ln n_j)$ (even); choosing $n_j=2^{2^j}$ gives $\sum_j\Delta_j/(n_j\operatorname{polylog})<\infty$ since $\Delta_j\le n_j^{o(1)}$. (P)
\item \emph{Maier-matrix irregularities, GPY positive-proportion-$\eta$ small gaps, LDG distributional statements:} accommodated by tuning $\xi$-statistics and the sprinkle profile. [HEURISTIC]
\end{itemize}
Every currently published unconditional theorem about primes is an upper-bound/existence statement compatible with these choices.

\subsection{Conclusion of conditional independence}

There is a self-consistent world $W(\theta,\rho)$, indistinguishable at the level of all proven prime statistics, in which $\sum_k\tilde b_k$ diverges like $\frac{1-\theta}{2}\ln\ln K$. Therefore \textbf{no currently unconditional theorem distinguishes our primes from a world where the series diverges.} The missing information is exactly \emph{parity-resolved local gap correlation} --- how $g_{2k-1}$ distributes around $p_{2k-1}/(2k-1)$ versus $g_{2k}$ around $p_{2k}/(2k)$. This is supplied by Hardy--Littlewood: Gallagher's calculation (HL implies Poissonian counts in $(x,x+\lambda\log x]$, mean parity $e^{-2\lambda}$), quantified via Kuperberg's singular-series moment estimates and the Banks--Ford--Tao random sifted model, yielding Tao's bound $\bigl|\sum_{n\le x}(-1)^{\pi(n)}\bigr|\ll x/(\log\log x)^{1.1}$ \cite{Ta23}.

\section{Numerical evidence}\label{sec:num}

Methodology \textnormal{[NUMERICAL]}: exact fixed-point arithmetic at scale $M=10^{60}$, round-half-up per term, rigorous accumulated rounding bound $\le N/(2M)$; validated against \texttt{fractions.Fraction} at $N=5000$ (absolute difference $1.581\times10^{-59}$); independent pass at $M'=10^{50}$ agrees to 40 significant digits; sieve self-checks $p_{10^6}=15485863$, $p_{10^7}=179424673$.

\begin{table}[h]
\centering\small
\begin{tabular}{@{}rlrr@{}}
\toprule
$N$ & $S_N=\sum_{n\le N}(-1)^n n/p_n$ (verbatim fixed-point value) & increment $\Delta_N=S_N-S_{N/2}$ & err.\ bound\\
\midrule
$10^3$ & \texttt{0.0118688917373843595188635122441718349053} & $-8.081\times10^{-3}$ & $5\times10^{-58}$\\
$10^4$ & \texttt{-0.0032901791643098451810940216336951591445} & $-3.068\times10^{-3}$ & $5\times10^{-57}$\\
$10^5$ & \texttt{-0.0135356832245282824221422477921787955820} & $-2.507\times10^{-3}$ & $5\times10^{-56}$\\
$10^6$ & \texttt{-0.0198592162984378367239838621029168996347} & $-1.611\times10^{-3}$ & $5\times10^{-55}$\\
$10^7$ & \texttt{-0.0242870341777922639490327584168847121429} & $-1.202\times10^{-3}$ & $5\times10^{-54}$\\
\bottomrule
\end{tabular}
\caption{Partial sums at exact fixed-point precision, half-decade increments, accumulated rounding bounds.}
\label{tab:S}
\end{table}

Interpretation:
\begin{itemize}\itemsep2pt
\item \emph{Increments decay steadily} ($-8.08\to-1.20$ in units of $10^{-3}$ per half-decade); no drift signal. Consistent with the slow conditional rate $O((\log\log x)^{-0.1})$ of \cite{Ta23} and with the putative limit $\approx-0.052161$ estimated there: $S_{10^7}=-0.0243$ is still far above it, as expected for logarithmic-rate convergence.
\item \emph{Term scale} $a_N=N/p_N$: $0.1263$ ($10^3$) $\to$ $0.0557$ ($10^7$).
\item \emph{Non-monotonicity empirical:} $\#\{n\le10^7: a_{n+1}<a_n\}/N=0.4213$ (58\% of steps increase) --- Leibniz fails empirically too, matching Fact~\ref{fac:up}.
\item \emph{Total variation (exact):} $TV(10^6)=1.4346694327924257960927888$, $TV(10^7)=1.5262825169363214597427909$; $TV/\ln\ln N\approx0.5464\to0.5490$, growing without apparent bound --- BV fails, matching Proposition~\ref{prop:tv}.
\item \emph{Pair sums} $B_K=S_{2K}$: $\max_{K\le5\cdot10^6}|B_K|=\tfrac16$ exactly at $K=1$ (transient first pair $\tfrac23-\tfrac12$); thereafter bounded well below. Paired variation $TV_{\mathrm{pair}}(10^7)=0.7534934775613393114390976$; cancellation ratio $\max|B_K|/TV_{\mathrm{pair}}(10^7)=0.2212$ --- massive sustained cancellation, as required for convergence.
\item \emph{Companion Erd\H{o}s series} $C_N=\sum_{n\le N}(-1)^n/(ng_n)$: $C(10^6)=-0.9697061021976592$, $C(10^7)=-0.9698267399607785$, half-decade drifts $\le2.87\times10^{-5}$ --- converging. (Open unconditionally and even under Conjecture~1.3 of \cite{Ta23}; Section~\ref{sec:sib}.)
\item \emph{Divergence certificate for the sibling} $\sum(-1)^n/g_n$: $\#\{n\le10^7:\ g_n=2\}=738597$ twin steps, and Maynard/Polymath give $g_n\le246$ infinitely often, so terms $\ge1/246$ persist along an infinite sequence.
\item \emph{Pre-registered divergence criteria:} divergence would manifest as ($\alpha$) non-vanishing amplitude of $\Delta_N=S_N-S_{N/2}$, or ($\beta$) sustained $\ln\ln$-drift of $B_K$ (cf.\ the $W(\theta)$ magnitude $\frac{|1-\theta|}{2}\ln\ln K$). Neither is observed through $10^7$.
\end{itemize}

\section{Sibling problems}\label{sec:sib}

\paragraph{D.1: $\sum(-1)^n/g_n$ diverges --- unconditional.}
Maynard/Polymath8b \cite{Ma15,Po14}: $g_n\le246$ infinitely often, so $|(-1)^n/g_n|=1/g_n\ge1/246$ along an infinite sequence; terms do not tend to zero, hence divergence (ordinary and alternating alike). Also recorded in \cite[\S5]{Ta23}. Contrast Weisenberg's observation [per erdosproblems.com page; UNVERIFIED-citation]: under HL $k$-tuples the partial sums are unbounded in at least one direction --- the sibling diverges robustly, while the weighted original resists proof.

\paragraph{D.2: weighted variants.}
Define
\[
U_c:=\sum_n\frac{1}{n\,g_n\,(\ln\ln n)^c}\ \ (\text{no sign}),\qquad
V_c:=\sum_n\frac{(-1)^n}{n\,g_n\,(\ln\ln n)^c}.
\]
\emph{(a) $U_c<\infty$ for $c>2$ (absolute-convergence theorem; reproduced).} Input (upper-bound-sieve class; cf.\ ``sparsity of very small gaps'' $\#\{p_n\le x:\ g_n\le h\}\ll\min(h/\ln x,1)\pi(x)$, \cite[Thm~2]{GPY11}; the sharp dyadic form is attributed on the problem page to a Selberg-sieve argument of Sawhney [UNVERIFIED-citation]):
$\#\{n\in[Y,2Y]:\ g_n\in[\varepsilon\ln n,2\varepsilon\ln n)\}\ll\varepsilon Y.$
Each term of the $\varepsilon$-class in the block satisfies $1/(ng_n(\ln\ln n)^c)\le\frac{2}{Y\cdot\varepsilon\lambda\cdot(\ln\lambda)^c}$ with $\lambda=\ln Y$, so the class contributes $\ll\frac{2}{\ln Y(\ln\ln Y)^c}$, independent of $\varepsilon$; with $J\approx\log_2\lambda$ active classes (truncate at $\varepsilon\ge1/\lambda$), the block total is
$\ll\frac{(\ln\ln Y)^{1-c}}{\ln Y}$
(the large-$g$ stratum is negligible by Markov counting $\#\{g\ge G\}\ll Y/G$ from $\sum g_n\asymp Y$). Summing over $Y=2^m$: $\sum_m\frac{(\ln m)^{1-c}}{m}$ converges $\iff c>2$. Matches Erd\H{o}s--Nathanson [per page; UNVERIFIED-citation] and Sawhney's sieve proof as relayed by the page.

\emph{(b) $U_c=\infty$ for $c\le1$ --- unconditional.} By GPY Theorem~1 \cite{GPY11}, $A_\eta=\{n:g_n\le\eta\ln p_n\}$ has lower density $\delta_\eta>0$ for every fixed $\eta>0$; on $A_\eta$, $\frac{1}{ng_n(\ln\ln n)^c}\ge\frac{1}{n\cdot\eta\ln n\cdot(\ln\ln n)^c}$, and the Window Lemma (Lemma~\ref{lem:window}) applied to $A_\eta$ gives
\[
\sum_{n\le N,\,n\in A_\eta}\frac{1}{ng_n(\ln\ln n)^c}
\gg_\eta\ \sum_{\text{windows }j}\frac{1}{\ln x_j(\ln\ln x_j)^c}
\asymp\sum_j\frac{1}{j(\ln j)^c}=\infty\quad\iff\quad c\le1 .
\]
(This strengthens the naive typical-gap heuristic $\sum\frac{1}{n\ln n(\ln\ln n)^c}$, which also diverges exactly for $c\le1$.)

\emph{(c) Regime $1<c\le2$ and the alternating sibling.} $V_c$ converges absolutely for $c>2$ by (a). At the endpoint, \cite[\S5]{Ta23} sketches (argument due to Ford) that $V_2$ \emph{diverges} assuming Conjecture~1.3 --- refuting, conditionally, the Er98 conjecture ``for all $c>0$'' at $c=2$. For $0\le c<2$ convergence of $V_c$ remains open even under Conjecture~1.3 \cite{Ta23}; probabilistic heuristics (random signs, Khintchine) support convergence of $V_0=\sum(-1)^n/(ng_n)$ --- the numerically convergent companion of Section~\ref{sec:num} --- and suggest more generally $\sum(-1)^n/(n^\theta g_n)$ convergent for $\theta>\tfrac12$ \cite{Ta23}.

\begin{remark}[erratum vs.\ the tasking sheet]
The sheet's claim ``absolute convergence FAILS for every $c$'' (via $\sum\frac{1}{ng_n(\ln\ln n)^c}\asymp\sum\frac{1}{n(\ln\ln n)^c}$) is incorrect as stated: the proposed comparison drops the essential $1/\ln n$ factor ($\mathbb E[1/g]\asymp\ln\ln n/\ln n$, not $\asymp1$; the exponential-law constant $\mathbb E|g-\ln n|\asymp\ln n$ concerns $|g-\ln n|$, not $1/g$), and $\sum\frac{1}{n\ln n(\ln\ln n)^c}$ itself converges for $c>1$. Correct picture: $U_c=\infty$ for $c\le1$ (proved), $U_c<\infty$ for $c>2$ (Erd\H{o}s--Nathanson/Sawhney), borderline resolved only conditionally at $c=2$.
\end{remark}

\paragraph{D.3: why the siblings bracket the original.}
Removing the weight ($1/g_n$ instead of $n/p_n$-type decay) destroys term-to-zero and gives unconditional divergence (D.1); keeping a stronger weight ($1/(ng_n)$) produces a series plausibly convergent but beyond even HL-strong at $c=0$. The original sits at the critical weighting $a_n\asymp1/\ln n$ where absolute convergence fails, monotonicity fails, BV fails --- and only parity-resolved information can decide.

\section{Boxed obstruction}\label{sec:obs}

\noindent\fbox{\parbox{\dimexpr\linewidth-2\fboxsep-2\fboxrule\relax}{
\textbf{Obstruction (the deliverable).}\enspace
The series $\sum_{n\ge1}(-1)^n n/p_n$ sits exactly in the dead zone of alternating-series theory: its terms $a_n=n/p_n$ decay only like $1/\ln n$, far too slowly for absolute convergence ($\sum1/\ln n$ diverges); the bounded-gaps theorems of Zhang--Maynard force $a_{n+1}>a_n$ infinitely often (indeed $a_{n+1}<a_n\iff ng_n>p_n$, and $g_n\le246\ll p_n/n$ i.o.), so Leibniz's monotonicity hypothesis is provably false in both directions; the exact step identity $a_{n+1}-a_n=(p_n-ng_n)/(p_np_{n+1})$, evaluated on the GPY positive-density set of gaps $\le\eta\ln p_n$, gives total variation $\gg\sum1/(n\ln n)\sim\ln\ln N\to\infty$ (observed: $TV(N)\approx0.55\ln\ln N$), killing the Dirichlet--BV criterion; and after subtracting the smooth Cipolla model $f(n)=n/E(n)$ --- whose alternating sum converges --- the residual $\sum(-1)^n\delta_n$, $\delta_n=-n(p_n-E(n))/(p_nE(n))$, has envelope $|\delta_n|\ll e^{-c\sqrt{\ln n}}$ (RH would merely sharpen it to $n^{-1/2+o(1)}$, which by itself decides nothing: $(-1)^nn^{-1/2}/\ln n$ meets that envelope with divergent alternating sum). All unconditional error-term technology is parity-blind and monotonicity-free: it bounds $|p_n-E(n)|$ but says nothing about the odd/even correlation that drives the exact pair terms $b_k=\bigl(p_{2k-1}-(2k-1)g_{2k-1}\bigr)/(p_{2k-1}p_{2k})$, on which convergence is exactly pinned by $S_{2K}=\sum_{k\le K}b_k$. The counterexample world $W(\theta,\rho)$ --- parity-biased bulk gaps whose pair-means $2\Lambda_k$, $\Lambda_k=\ln(2k)+\ln\ln(2k)$, telescope to $\tilde p_N=N(\ln N+\ln\ln N-1)+O(N/\ln N)=N\cdot$(full Cipolla main terms)$+o(N)$, compensated bounded-gap sprinkle, sparse huge gaps --- satisfies every published unconditional prime statistic while its pair series diverges like $\frac{1-\theta}{2}\ln\ln K$ (the systematic non-parity pieces of $\tilde b_k$ form convergent series $\sum O(1/(k\ln^2k))$); hence resolution demands genuinely new, \emph{parity-resolved local prime statistics}. Exact missing bridge: \emph{``uniform-in-$\lambda$ parity-resolved asymptotics for $\#\{k\le K:\ g_{2k-1}\le\lambda p_{2k-1}/(2k-1)\}$ (equivalently parity-equidistribution of $\pi(m)$ in log-length intervals / a power-saving $\sum_{n\le x}(-1)^{\pi(n)}\ll x/(\log\log x)^{1.1}$) $\Rightarrow\sum_kb_k$ converges $\Rightarrow$ EP-15 convergent''} --- supplied today only by the Quantitative Hardy--Littlewood prime-tuples conjecture via Tao's random-sifted-model proof; nothing weaker is known, and nothing weaker can suffice by Section~\ref{sec:r5}.
}}

\section{Exact missing implication}\label{sec:chain}

\[
\boxed{
\begin{gathered}
\text{uniform-in-}\lambda\text{ parity-resolved asymptotics}\\
\#\{k\le K:\ g_{2k-1}\le\lambda\, p_{2k-1}/(2k-1)\}=\mu(\lambda)K+o(K)\ \ (\lambda\ \text{in a range}),\\
\Updownarrow\ \text{(equiv.)}\\
\text{equidistribution of the parity of }\pi(m)\text{ in intervals }(x,x+\lambda\log x],\ \text{i.e.}\ \sum_{n\le x}(-1)^{\pi(n)}\ll x/(\log\log x)^{1.1},\\
\Downarrow\\
\text{the signed pair series }\sum_kb_k\text{ converges}\\
\Downarrow\\
\textstyle\sum(-1)^n n/p_n\ \text{converges.}
\end{gathered}}
\]
Tao \cite{Ta23} establishes the needed consequence of this chain modulo his Conjecture~1.3 (Quantitative HL prime tuples); nothing weaker is known to supply it --- and by Section~\ref{sec:r5}, anything compatible with $W(\theta\ne1,\rho)$ provably cannot.

\section*{Errata recorded against the original tasking sheet}
\addcontentsline{toc}{section}{Errata recorded against the original tasking sheet}
\begin{enumerate}\itemsep2pt
\item \emph{R4 sign:} correct chain $2kp_{2k-1}-(2k-1)p_{2k}=p_{2k-1}-(2k-1)g_{2k-1}$; hence $b_k>0\iff g_{2k-1}<p_{2k-1}/(2k-1)$ (sheet had the inverted direction; verified exactly for $k\le50{,}000$; $b_1=1/6$ sanity passes).
\item \emph{R5 probe:} executed script summed $(\theta-1)\sum1/(2kL_k^2)$ whose total converges; modeled pair term $b_k=(L_k-g)/(2kL_k^2)$ with $g=\theta L_k$ gives $(1-\theta)/(2kL_k)$, genuinely diverging like $\frac{1-\theta}{2}\ln\ln K$; sign also corrected.
\item \emph{D.2:} the sheet's ``absolute divergence for every $c$'' is incorrect as stated; correct regimes as in Section~\ref{sec:sib}.
\end{enumerate}

\begin{thebibliography}{99}\small
\raggedright

\bibitem{Ta23} T.~Tao, \emph{The convergence of an alternating series of Erd\H{o}s, assuming the Hardy--Littlewood prime tuples conjecture}, arXiv:2308.07205 (2023).
\bibitem{Ta24} T.~Tao, \emph{The convergence of an alternating series of Erd\H{o}s, assuming the Hardy--Littlewood prime tuples conjecture}, Comm.\ Amer.\ Math.\ Soc.\ \textbf{4} (2024), no.~3, 80--96.
\bibitem{Er98} P.~Erd\H{o}s, \emph{Some of my new and almost new problems and questions on combinatorial number theory, prime and additive number theory}, Number theory (Eger, 1996), de Gruyter, Berlin, 1998, pp.~169--180. [Problem E7 at p.~203 referenced via \cite{Ta23}; sibling conjectures as relayed by the problem pages.]
\bibitem{Zh14} Y.~Zhang, \emph{Bounded gaps between primes}, Ann.\ of Math.\ (2) \textbf{179} (2014), 1121--1174.
\bibitem{Ma15} J.~Maynard, \emph{Small gaps between primes}, Ann.\ of Math.\ (2) \textbf{181} (2015), 383--413.
\bibitem{Po14} D.\,H.\,J.~Polymath, \emph{Variants of the Selberg sieve, and bounded intervals containing many primes}, J.\ Austral.\ Math.\ Soc.\ \textbf{97} (2014), 321--347. [$\liminf g_n\le246$.]
\bibitem{GPY11} D.~A.~Goldston, J.~Pintz, C.~Y.~Y{\i}ld{\i}r{\i}m, \emph{Positive proportion of small gaps between consecutive primes}, Publ.\ Math.\ Debrecen \textbf{79} (2011), 433--444; arXiv:1103.3986. [Theorem~1: $P(x,\eta)\gg_\eta1$; Theorem~2: sparsity of very small gaps. See also \emph{Primes in tuples IV}, Acta Arith.\ \textbf{160} (2013).]
\bibitem{We31} \O{}.~Westzynthius, \emph{\"{U}ber die Verteilung der Zahlen, die zu den $n$ ersten Primzahlen teilerfremd sind}, Comment.\ Phys.-Math.\ Soc.\ Sci.\ Fennicae \textbf{5} (1931), no.~25, 1--37.
\bibitem{Ci02} M.~Cipolla, \emph{La determinazione asintotica dell'$n$-mo numero primo}, Matematiche (Catania) \textbf{8} (1902), 132--146.
\bibitem{vK01} H.~von Koch, \emph{Sur la distribution des nombres premiers}, Acta Math.\ \textbf{24} (1901), 159--182.
\bibitem{DVLP} Ch.~J.~de la Vall\'ee Poussin, \emph{Recherches analytiques sur la th\'eorie des nombres premiers}, Ann.\ Soc.\ Sci.\ Bruxelles \textbf{20} (1896).
\bibitem{Ga76} P.~X.~Gallagher, \emph{On the distribution of primes in short intervals}, Israel J.\ Math.\ \textbf{23} (1976), 193--198.
\bibitem{BFT} W.~D.~Banks, K.~Ford, T.~Tao, paper introducing the random sifted model of the primes (referenced as [1]/[2] in \cite{Ta23}). [UNVERIFIED-citation: exact coordinates not checked.]
\bibitem{Ku} V.~Kuperberg, quantitative estimates for mean values of the singular series (referenced as [11] in \cite{Ta23}). [UNVERIFIED-citation: exact coordinates not checked.]
\bibitem{Bl} B.~Bloom, \texttt{erdosproblems.com/15} (page relaying EP-15 status, Er98 conjectures, Weisenberg and Sawhney remarks). [Page facts as given; Sawhney/Weisenberg/Erd\H{o}s--Nathanson citation details UNVERIFIED-citation.]
\bibitem{UM} UnsolvedMath, entry EP-15. [Status Open; facts relayed in the tasking confirmed against \cite{Ta23} and \cite{GPY11}.]
\bibitem{MO} Anonymous (Said), unpublished observation, \texttt{mathoverflow.net/questions/313999} (equivalence recorded in \cite[Tao, footnote~3]{Ta23}).

\end{thebibliography}

\end{document}
